\documentclass[11pt]{article}
\usepackage[a4paper,margin=1in]{geometry}
\usepackage[T1]{fontenc}
\usepackage{lmodern}
\usepackage{microtype}
\usepackage{amsmath,amssymb,mathtools}
\usepackage{booktabs,array}
\usepackage{graphicx}
\usepackage{hyperref}
\usepackage{xcolor}
\usepackage{enumitem}
\hypersetup{colorlinks=true,linkcolor=blue!55!black,urlcolor=blue!55!black}
\setlist{itemsep=0.3em,topsep=0.5em}
\setlength{\parskip}{0.45em}
\setlength{\parindent}{0pt}
\title{Cosmological Expansion and Dark Matter as Emergent\\
Informational Phenomena}
\author{Juha Meskanen}
\date{July 2026 -- implementation-aligned revision}

\begin{document}
\maketitle

\begin{abstract}
This paper documents the cosmological model as it is implemented in the
current simulation code.  The universe is represented by a finite static
informational structure of $n$ binary degrees of freedom.  The simulation
uses a raw internal ordering coordinate $T$ and the normalized coordinate
$\tau=T/n$.  Its substrate layer follows the continuum Ehrenfest mean
\begin{equation*}
 p(\tau)=\frac12\left(1-e^{-2\tau}\right),
 \qquad
 k(T)=\operatorname{round}\!\bigl(np(T)\bigr),
\end{equation*}
and defines the unfolded informational pool by
$U(T)=\log_2\binom{n}{k(T)}$.  A width-$w$ window is classified with the
exact finite-population hypergeometric distribution, conditional on the
implemented integer count $k(T)$.  The three classes used by the code are
$J=0$ (falling), $0<J<\lceil w/2\rceil$ (structure-forming), and
$J\geq\lceil w/2\rceil$ (promoted).

Each scale has an internal tick count $q_w(T)=\lfloor T/w\rfloor$ and is
evaluated at the retarded raw coordinate $T_w=wq_w(T)$.  Recursively
promoted pools form a nested classical hierarchy.  At every level the
available pool is partitioned into fabric, structure, promoted, and pending
sectors, and the resulting telescoping information ledger is exact to
floating-point precision.  The code defines matter bits by
$M_i(T)=C_i(T)\eta(T)^q$, where $C_i$ is the structure sector,
$\eta=e^{-2T/n}$, and the exponent $q$ is an explicit non-derived modeling
choice whose current default is $q=1$.  The implemented relational size
measure is
\begin{equation*}
 R_D(T)=\frac{U(T)-\sum_i L_i(T)-\sum_i M_i(T)}{n},
\end{equation*}
where $L_i$ is the pending multi-clock backlog.  Composite entity counts
$N_i=M_i/w_i$ are used for worldlines and abundance plots but are not added
back into $R_D$; in the current implementation matter bits are unavailable
as free spacetime fabric.

The default run displays three descriptive regimes in internal coordinate
time: rapid growth of the unfolded pool, a matter-loaded and clock-retarded
regime in which growth is suppressed and local contractions occur, and, for
$q>0$, recovery toward the no-matter entropy envelope as the matter weight
decays.  These are properties of the implemented finite-information model,
not yet a derivation of an FLRW metric, physical acceleration, or a calibrated
cosmic time.  A companion $\psi$-layer implementation represents fixed-count
sectors by Dicke states, uses the same hypergeometric window statistics, and
supports both genuinely nested cascades and parallel coexisting species.
The paper separates exact identities, implemented approximations, modeling
choices, and future physical identifications.
\end{abstract}

\tableofcontents
\new page

\section{Purpose and Scope}
\label{sec:scope}

The foundational ontology of the framework is a finite, static, timeless
informational universe.  Time is internal to an observer: it is an ordering
of distinguishable observer-compatible states, not an external parameter in
which the total universe evolves.  The simulation coordinate $T$ used below
is therefore an ordinal bookkeeping variable for elementary informational
transitions.  Running the Python program in laboratory time is a method for
studying the static relational model; processor time is not identified with
cosmic time.

The purpose of this paper is narrower than a complete cosmological theory.
It records the present mathematical content of the implementation in
\path{multiclock.py}, its classical driver \path{cosmic_d.py}, and the
spectral companion implemented by \path{dicke_layer.py},
\path{dicke_cascade.py}, and \path{cosmic_psi.py}.  In particular, the
paper distinguishes four levels of statement:
\begin{enumerate}
  \item exact finite combinatorial identities;
  \item the mean-shell and retarded-clock approximations actually used by
        the code;
  \item explicit modeling choices, such as the matter exponent and window
        compositions;
  \item proposed physical interpretations, such as spacetime curvature,
        dark matter, and cosmological expansion.
\end{enumerate}

The code is treated as the authoritative specification of the present
statistical model.  This does not make every comment in the source a theorem.
Where the implementation uses a continuum approximation to a discrete chain,
the paper states that qualification explicitly.

\section{Finite Substrate and Internal Ordering}
\label{sec:substrate}

\subsection{The discrete Ehrenfest chain and the implemented mean curve}

Let $K_T$ denote the number of ones after $T$ elementary bit flips in a
length-$n$ bitstring.  One uniformly selected bit is toggled at each discrete
tick.  Conditional on $K_T=k$,
\begin{equation}
 \Pr(k\to k+1)=\frac{n-k}{n},
 \qquad
 \Pr(k\to k-1)=\frac{k}{n}.
 \label{eq:birthdeath}
\end{equation}
Hence
\begin{equation}
 \mathbb{E}[K_{T+1}-K_T\mid K_T=k]=1-\frac{2k}{n}.
 \label{eq:drift}
\end{equation}
Starting from $K_0=0$, the exact discrete expectation is
\begin{equation}
 \frac{\mathbb{E}K_T}{n}
 =\frac12\left[1-\left(1-\frac2n\right)^T\right].
 \label{eq:discrete-mean}
\end{equation}

The current implementation uses the continuum, or independent-Poisson,
form
\begin{equation}
 p(T)=\frac12\left(1-e^{-2T/n}\right)
     =\frac12\left(1-e^{-2\tau}\right),
 \qquad \tau=\frac{T}{n}.
 \label{eq:p-code}
\end{equation}
The coefficient $2$ is derived from the Ehrenfest drift and is not fitted.
For the one-flip-per-tick discrete chain, equation~\eqref{eq:p-code} is the
large-$n$ approximation to equation~\eqref{eq:discrete-mean}; it is exact for
a corresponding continuous-time process with independent Poisson flip
clocks.  The code then uses the mean shell
\begin{equation}
 k(T)=\operatorname{clip}\!\left(
       \operatorname{round}[np(T)],0,n\right).
 \label{eq:k-code}
\end{equation}
Thus the implementation does not propagate the complete probability
distribution of $K_T$.  It evaluates finite-window statistics conditional on
the rounded mean count.

\subsection{Unfolded informational pool}

The code defines the substrate's unfolded pool by the fixed-count
combinatorial entropy
\begin{equation}
 U(T)=\log_2\binom{n}{k(T)}.
 \label{eq:unfolded}
\end{equation}
The plotted quantity named \texttt{entropy} is
\begin{equation}
 E_n(T)=\frac{U(T)}{n}.
 \label{eq:entropy-normalized}
\end{equation}
For finite even $n$, the maximum is
$\log_2\binom{n}{n/2}/n<1$, approaching one as $n\to\infty$.
This $U(T)$ is the budget passed into the first scale of the hierarchy.
It is not the same quantity as the contrast coordinate $2p(T)$, although
both increase toward saturation on the relaxation-from-zero branch.

The departure from equilibrium used in the matter ansatz is
\begin{equation}
 \eta(T)=1-2p(T)=e^{-2T/n}.
 \label{eq:eta}
\end{equation}

\subsection{Relation to the theoretical $n/2$ contrast bound}

An internal observer who constructs distance from complementary zero-one
distinctions can form at most
\begin{equation}
 Q_{\max}=\left\lfloor\frac n2\right\rfloor
 \label{eq:n2bound}
\end{equation}
independent elementary contrasts.  This bound is independent of the
dimension or topology assigned by a later decoder.  The current simulation
does not explicitly pair zeros and ones.  Instead it uses the entropy pool
$U(T)$ as an addressable-fabric proxy and normalizes by $n$.  Consequently,
$R_D$ below is the size measure implemented in the code, not yet the final
contrast-count metric.  A future metric decoder may rescale or replace this
proxy using equation~\eqref{eq:n2bound}.

\section{Exact Conditional Window Statistics}
\label{sec:windows}

For a width-$w$ window in a global shell containing exactly $k$ ones, let
$J$ be the number of ones in the window.  The conditional distribution is
hypergeometric:
\begin{equation}
 h_w(j\mid n,k)
 =\Pr(J=j\mid K=k)
 =\frac{\binom{k}{j}\binom{n-k}{w-j}}{\binom{n}{w}}.
 \label{eq:hypergeom}
\end{equation}
Equation~\eqref{eq:hypergeom} is exact for a uniformly sampled fixed-$k$
shell.  In the time-series code it is evaluated at the mean-shell value in
equation~\eqref{eq:k-code}.

Define the threshold
\begin{equation}
 j_*(w)=\left\lceil\frac w2\right\rceil.
 \label{eq:threshold}
\end{equation}
The three code-level families are
\begin{align}
 f^{(0)}_w(T)
   &=h_w(0\mid n,k(T)),
   \label{eq:fall}\\
 f^{(b)}_w(T)
   &=\sum_{j=1}^{j_*(w)-1}h_w(j\mid n,k(T)),
   \label{eq:bump}\\
 f^{(r)}_w(T)
   &=\sum_{j=j_*(w)}^{w}h_w(j\mid n,k(T)).
   \label{eq:rise}
\end{align}
Here $f^{(0)}_w$ is the falling all-zero class, $f^{(b)}_w$ is the
structure-forming or bump class, and $f^{(r)}_w$ is the promoted or rising
class.
They obey
\begin{equation}
 f^{(0)}_w+f^{(b)}_w+f^{(r)}_w=1.
 \label{eq:family-partition}
\end{equation}
The explicit separation of $J=0$ is important.  A composition with no ones
is monotonically falling and is not an interior hump.

For comparison, one exact ordered pattern with $a$ required ones and $b$
required zeros has, in the independent-bit limit,
\begin{equation}
 P_{a,b}(p)=p^a(1-p)^b.
 \label{eq:pattern}
\end{equation}
For $0<a<b$ it peaks at $p_*=a/(a+b)<1/2$.  The classical cosmological
pipeline does not select one fixed $(a,b)$ pattern; it aggregates all
window compositions in equations~\eqref{eq:fall}--\eqref{eq:rise}.

\section{Multi-Clock Hierarchy}
\label{sec:multiclock}

\subsection{Scale-dependent ticks}

A width-$w$ composite requires $w$ elementary transitions per internal
update in the present clock model.  Its own tick count is
\begin{equation}
 q_w(T)=\left\lfloor\frac{T}{w}\right\rfloor,
 \label{eq:own-ticks}
\end{equation}
and its state is evaluated at the retarded raw coordinate
\begin{equation}
 T_w=wq_w(T)=w\left\lfloor\frac{T}{w}\right\rfloor.
 \label{eq:retard}
\end{equation}
The quantity $q_w$ advances, on average, at rate $1/w$ relative to the raw
ordering coordinate.  The code stores this as a nominal lapse
$\ell_w=1/w$.  Equation~\eqref{eq:retard} is a sample-and-hold rule: the
scale's state changes only when another block of $w$ raw transitions has
accumulated.

\subsection{Recursive level state}

Consider $m$ levels with widths $w_0,\ldots,w_{m-1}$.  Let $A_i(T)$ be the
true pool presented to level $i$ at raw coordinate $T$, and let
$\widetilde A_i(T)$ be the pool visible to that level at its retarded
coordinate.  The first pool is
\begin{equation}
 A_0(T)=U(T),
 \qquad
 \widetilde A_0(T)=U(T_{w_0}).
 \label{eq:first-pool}
\end{equation}
At level $i$, define
\begin{align}
 L_i(T)&=A_i(T)-\widetilde A_i(T),
 \label{eq:pending}\\
 F_i(T)&=\widetilde A_i(T) f^{(0)}_{w_i}(T_{w_i}),
 \label{eq:fabric}\\
 C_i(T)&=\widetilde A_i(T) f^{(b)}_{w_i}(T_{w_i}),
 \label{eq:structure}\\
 G_i(T)&=\widetilde A_i(T) f^{(r)}_{w_i}(T_{w_i}).
 \label{eq:promoted}
\end{align}
These are, respectively, the pending backlog, fabric sector, raw structure
sector, and promoted sector.
The next level receives the promoted output,
\begin{equation}
 A_{i+1}(T)=G_i(T).
 \label{eq:recursion}
\end{equation}
In the implementation, evaluating $A_i(T_{w_i})$ recursively retards all
earlier levels consistently; it is not obtained by merely shifting the
last array.

Equations~\eqref{eq:family-partition} and \eqref{eq:pending} imply
\begin{equation}
 A_i=L_i+F_i+C_i+G_i.
 \label{eq:level-ledger}
\end{equation}
Using equation~\eqref{eq:recursion}, the hierarchy telescopes to
\begin{equation}
 \boxed{
 U(T)=\sum_{i=0}^{m-1}\bigl[L_i(T)+F_i(T)+C_i(T)\bigr]
      +G_{m-1}(T).}
 \label{eq:global-ledger}
\end{equation}
This is the conservation identity tested by
\texttt{conservation\_max\_error}.  It is an exact consequence of the
implemented partition, up to floating-point error.

The classical pipeline is therefore a genuinely nested formation hierarchy.
A later level is built only from the promoted sector of the preceding level.
This is appropriate when the levels represent successive stages of one
formation cascade.  It is not automatically appropriate for distinct
species that coexist without being nested; the spectral implementation
provides a separate parallel mode for that case.

\section{Matter Definition and Implemented Size Measure}
\label{sec:size}

\subsection{Matter weighting}

The raw structure sector $C_i$ is converted into the code's matter sector by
\begin{equation}
 M_i(T)=C_i(T)\,\eta(T)^q,
 \label{eq:matter}
\end{equation}
where $q$ is the command-line parameter \texttt{matter\_power}.  The current
default is $q=1$.  The total matter-bit allocation and the associated entity
count are
\begin{equation}
 M(T)=\sum_i M_i(T),
 \qquad
 N(T)=\sum_i\frac{M_i(T)}{w_i}.
 \label{eq:matter-total}
\end{equation}
The auto-selected marker called ``today'' in the figures is
\begin{equation}
 T_{\rm today}=\operatorname*{arg\,max}_T N(T).
 \label{eq:today}
\end{equation}
It is a plotting convention, not a derived physical-age calibration.

The exponent $q$ is not fixed by the window combinatorics.  At $q=0$, every
raw structure bit counts as matter and the equilibrium structure plateau
survives.  For $q>0$, equation~\eqref{eq:eta} suppresses matter toward late
internal times.  The model's late behavior therefore depends strongly on a
currently non-derived choice.

\subsection{Size measure used by \texttt{run\_simulation}}

Let
\begin{equation}
 L(T)=\sum_iL_i(T)
 \label{eq:pending-total}
\end{equation}
be the total multi-clock backlog.  The code defines
\begin{equation}
 \boxed{
 R_D(T)=\frac{U(T)-L(T)-M(T)}{n}.}
 \label{eq:size-code}
\end{equation}
This equation replaces the older manuscript formula in which composite
entity counts were added back into the size measure.  In the present code,
$N_i=M_i/w_i$ is used to display matter worldlines and abundances, but it is
not included in equation~\eqref{eq:size-code}.  A $w_i$-bit matter structure
therefore withdraws all $w_i$ participating bits from the free spacetime
size budget.

Using the exact ledger~\eqref{eq:global-ledger}, equation~\eqref{eq:size-code}
can also be written as
\begin{equation}
 nR_D
 =\sum_iF_i+\sum_i(C_i-M_i)+G_{m-1}.
 \label{eq:size-ledger}
\end{equation}
This identity exposes another current modeling convention.  When $q>0$,
only the fraction $\eta^q$ of the raw structure sector is called matter;
$C_i-M_i$ remains part of the size measure.  At $q=0$, all raw structure is
matter and that term vanishes.

\subsection{Informational interpretation}

The mechanism represented by equation~\eqref{eq:size-code} is simple.  No
new bits enter the static universe.  Bits that are available as independently
addressable fabric contribute to spatial resolution.  The same bits can be
organized into larger observer and matter microstructures.  Once allocated
to such a composite, they are no longer counted as free fabric.  Spatial
contraction and slower local clocks are therefore two readings of one finite
budget:
\begin{equation}
 \text{composite formation}
 \quad\Longrightarrow\quad
 \begin{cases}
  \text{fewer free spatial distinctions},\\
  \text{more elementary transitions per composite tick}.
 \end{cases}
 \label{eq:curvature-intuition}
\end{equation}
The present simulation is a homogeneous global shadow of this idea.  It does
not yet calculate a local metric, a Ricci tensor, or directional focusing.
Those require a local, direction-dependent version of the counting deficit.

\section{Spectral Companion}
\label{sec:spectral}

The $\psi$-layer companion represents a fixed global count by the Dicke state
\begin{equation}
 \lvert D_n^k\rangle
 =\frac{1}{\sqrt{\binom nk}}
  \sum_{x:\,|x|=k}\lvert x\rangle.
 \label{eq:dicke}
\end{equation}
For a width-$w$ subsystem, the probability of observing $j$ excitations is
again equation~\eqref{eq:hypergeom}.  The reduced-state entropy is therefore
\begin{equation}
 S_{\rm vN}(n,k,w)
 =-\sum_j h_w(j\mid n,k)\log_2h_w(j\mid n,k).
 \label{eq:svn}
\end{equation}
The code verifies the rank-one block structure underlying this shortcut by
direct state-vector diagonalization for small systems.

For a specified composition $(a,b)$ with $w=a+b$, the probability of the
whole composition class is
\begin{equation}
 P_{\rm class}(a,b\mid n,k)=h_w(a\mid n,k),
 \label{eq:classprob}
\end{equation}
whereas one exact ordered pattern has probability
\begin{equation}
 P_{\rm specific}(a,b\mid n,k)
 =\frac{h_w(a\mid n,k)}{\binom wa}.
 \label{eq:specificprob}
\end{equation}
The cosmological matter pipeline in \texttt{cosmic\_psi.py} deliberately
uses the class probability, because matter formation is associated with the
composition class rather than one chosen ordering.

The conservative persistence probability used by the spectral cascade is
\begin{equation}
 s(n,w)=\left(\frac{n-w}{n}\right)^w.
 \label{eq:survival}
\end{equation}
It is the probability that none of the $w$ bits is touched during the next
$w$ raw flips.  This literal-freeze criterion is exact for the stated event
but is a lower bound on more permissive composition-preserving persistence.

Two spectral organizations are implemented:
\begin{description}
 \item[Nested cascade.]  Each level acts on the exact leftover substrate of
 the preceding composition, and match and survival probabilities multiply
 across levels.  This models a genuine formation hierarchy.
 \item[Parallel species.]  Every level reads the same shared $(n,k(T))$ and
 receives only its own survival factor.  No cross-species multiplication is
 applied.  This is the appropriate implemented option for coexisting
 categories that are not nested inside one another.
\end{description}

For a level with $n_i$ available substrate bits, the spectral matter-bit
proxy is
\begin{equation}
 M_i^{\psi}(T)
 =\left\lfloor\frac{n_i}{w_i}\right\rfloor w_i
   P_i^{\rm persistent}(T),
 \label{eq:psi-matter}
\end{equation}
and the plotted spectral size proxy is
\begin{equation}
 R_{\psi}(T)
 =\max\left\{0,\frac{U(T)-\sum_iM_i^{\psi}(T)}{n}\right\}.
 \label{eq:psi-size}
\end{equation}
The current default composition
$a=\max(1,\lfloor w/3\rfloor)$, $b=w-a$ is explicitly a modeling choice,
not a particle-spectrum derivation.

The classical and spectral simulations are companion descriptions rather
than identical algorithms.  The classical code uses aggregated threshold
families, a nested multi-clock backlog, and the free exponent $q$.  The
spectral code uses selected composition classes, literal-freeze persistence,
and either nested or parallel organization.  Their common exact object is
the fixed-count hypergeometric window marginal.

\section{Default Numerical Experiment}
\label{sec:numerics}

\subsection{Protocol}

The default classical run uses the settings in
Table~\ref{tab:defaults}.
\begin{table}[ht]
\centering
\caption{Default parameters of \texttt{cosmic\_d.py}.}
\label{tab:defaults}
\begin{tabular}{lll}
\toprule
Quantity & Code value & Status \\
\midrule
Bit count & $n=184$ & illustrative finite system \\
Widths & $(6,12,20)$ & placeholders \\
Grid points & $3000$ & numerical display choice \\
Maximum raw coordinate & $T_{\max}=n\ln n$ & display-range choice \\
Ehrenfest rate & $2$ & derived continuum rate \\
Matter exponent & $q=1$ & non-derived current default \\
``Today'' & peak of $N(T)$ & plotting convention \\
Physical-year labels & log interpolation & plotting only \\
\bottomrule
\end{tabular}
\end{table}

The optional argument \texttt{--t\_bf\_max} is interpreted in units of $n$
and then multiplied by $n$.  If it is omitted, the endpoint is $n\ln n$.
The year-axis function maps the Planck-time label and $13.8$ billion years
log-uniformly onto the internal coordinate, with the ``now'' label placed at
equation~\eqref{eq:today}.  This map is not used in any model equation.

\subsection{Reproduced diagnostics}

Running the current default implementation gives
\begin{align*}
 \max_T|\text{ledger error}|&=2.84\times10^{-14}\ \text{bits},\\
 \frac{p(T_{\max})}{1/2}&=0.999970,\\
 T_{\rm today}&=60.1517,
 \qquad \tau_{\rm today}=0.326912,\\
 \max_T N(T)&=9.02814\ \text{entities},\\
 E_n(T_{\max})&=0.9777745,\\
 R_D(T_{\max})&=0.9777530.
\end{align*}
The distinct retarded raw states visited by the three scales over the
$3000$-point display grid are $160$, $80$, and $48$, respectively.  Their
maximum pending fractions are approximately $0.1931$, $0.02885$, and
$0.04972$ of the total bit budget.

Figure~\ref{fig:default} shows the four quantities that directly enter the
implemented size equation.  The step structure is not numerical noise; it is
produced by the retarded coordinates in equation~\eqref{eq:retard}.
\begin{figure}[ht]
\centering
\includegraphics[width=0.96\textwidth]{figures/cosmological-model-current-default.png}
\caption{Current default implementation in normalized internal coordinate
$\tau=T/n$.  The no-matter entropy pool $U/n$ rises rapidly.  Matter bits and
the pending multi-clock backlog suppress the implemented size measure
$R_D$.  The dashed line is the code's auto-selected ``now'', defined as the
peak total entity count.}
\label{fig:default}
\end{figure}

\subsection{Dependence on the matter exponent}

The asymptotic behavior is controlled by the non-derived exponent $q$.
Table~\ref{tab:qscan} reproduces a scan with all other default settings held
fixed.
\begin{table}[ht]
\centering
\caption{Sensitivity to the matter exponent $q$.  The relic column is the
final total entity count divided by its peak value.}
\label{tab:qscan}
\begin{tabular}{ccc}
\toprule
$q$ & $N(T_{\max})/N_{\max}$ & $R_D(T_{\max})$ \\
\midrule
$0$   & $7.343\times10^{-1}$ & $0.24984$ \\
$0.5$ & $6.440\times10^{-3}$ & $0.97382$ \\
$1$   & $4.930\times10^{-5}$ & $0.97775$ \\
$2$   & $2.400\times10^{-9}$ & $0.97777$ \\
\bottomrule
\end{tabular}
\end{table}

\begin{figure}[ht]
\centering
\includegraphics[width=0.96\textwidth]{figures/cosmological-model-matter-power.png}
\caption{The implemented size measure for several values of the free matter
exponent.  The $q=0$ curve retains the equilibrium structure plateau.  For
$q>0$, matter is increasingly suppressed by $\eta^q$ and $R_D$ returns
toward the no-matter entropy envelope.}
\label{fig:qscan}
\end{figure}

The late recovery of the default curve is therefore an implemented
consequence of the $q=1$ matter ansatz, not a theorem following from the bare
hypergeometric structure fraction.  Any cosmological interpretation of the
late regime must either derive $q$ or replace equation~\eqref{eq:matter} by a
persistence rule with no free exponent.

\subsection{Large-$n$ behavior}

At fixed normalized coordinate $\tau$ and fixed finite widths, the continuum
mean curve is independent of $n$, the normalized combinatorial entropy
approaches the binary entropy, and the hypergeometric marginal approaches a
binomial marginal.  The simulations consequently show a stable large-$n$
shape with finite-size and finite-width corrections.  This is approximate
convergence, not exact equality for all $n$.  In particular, the clock
backlog associated with a fixed width scales differently when $w/n$ is not
small.

\section{What the Implemented Profile Means}
\label{sec:profile}

The default output can be described in three internal-coordinate regimes.
The terminology here is deliberately kinematic rather than observational.

\paragraph{I. Entropy-pool growth.}
At small $T$, $U(T)$ grows rapidly from zero.  Before the matter and pending
terms become large, $R_D$ follows the same upward tendency.

\paragraph{II. Matter-loaded, clock-retarded growth.}
The structure fractions rise, matter is allocated from the same finite
budget, and scale-dependent updates create pending backlogs.  Consequently
$R_D<U/n$.  Each retarded scale produces plateaus and short contractions when
its local state updates.  The code therefore implements both a reduced
spatial budget and slower composite clocks from the same width parameter.

\paragraph{III. Late recovery for $q>0$.}
The order parameter $\eta^q$ suppresses matter, and pending backlogs vanish as
the hierarchy saturates.  The size measure approaches the no-matter entropy
envelope.  For $q=0$, this recovery is incomplete because the raw equilibrium
structure fraction remains classified as matter.

These regimes should not yet be called inflation, deceleration, or dark-energy
acceleration in the strict GR/FLRW sense.  Such statements require a physical
proper-time decoder, an operational metric, a redshift law, and derivatives
with respect to that proper time.  The implemented result is instead:
\begin{quote}
A fixed finite information budget, together with scale-dependent clocks and
matter allocation from the same substrate, generically suppresses the
observer's available spatial-resolution proxy relative to the no-matter
entropy envelope and can produce local contraction followed by recovery.
\end{quote}

\section{Dark Matter Interpretation}
\label{sec:darkmatter}

The code itself contains widths and compositions, not experimentally
identified particle names.  An assignment of one emergent class to dark
matter is therefore an interpretation layered on the statistical model.
Within the broader framework, the proposed mechanism is:
\begin{enumerate}
  \item an early, relatively simple fermionic microstructure class forms from
        the same finite fabric budget;
  \item its participating bits reduce the free spacetime capacity through
        equation~\eqref{eq:size-code}, so it gravitates by the same counting
        mechanism as visible matter;
  \item its compression residual is not in the photon class registered by
        visible-matter detectors, so it is electromagnetically dark;
  \item it may coexist with visible species, in which case the spectral
        parallel implementation is more appropriate than a forced nested
        cascade.
\end{enumerate}

This introduces no separate fundamental substrate, but it does define a
distinct emergent matter class.  The present simulation has not yet derived
its stability, velocity distribution, clustering, pressure, abundance, or
perturbation dynamics.  It therefore supplies an ontology and a formation
mechanism, not yet a complete cosmological dark-matter model.

\section{Connection to Curvature and General Relativity}
\label{sec:gr}

The global size equation is the homogeneous statistical shadow of a proposed
local counting mechanism.  At a point $x$, let a direction-dependent count
$\Delta N(x,k)$ describe the loss of independently usable fabric relations in
a small causal pencil oriented along a null vector $k^\mu$.  A continuum
Ricci limit would require a relation of the form
\begin{equation}
 \Delta N(x,k;A,\ell)
 =\gamma A\ell^2 T_{\mu\nu}(x)k^\mu k^\nu
 +o(A\ell^2),
 \label{eq:directional-count}
\end{equation}
combined with the geometric area-deficit expansion
\begin{equation}
 \frac{\Delta A}{A}
 =-\frac12\ell^2 R_{\mu\nu}(x)k^\mu k^\nu+o(\ell^2).
 \label{eq:ricci-target}
\end{equation}
The present cosmological code does not implement equations
\eqref{eq:directional-count}--\eqref{eq:ricci-target}.  Its contribution is
the finite-budget principle and the scale-dependent clock rule that a local
tensor construction would have to refine.  Demonstrating the Ricci tensor
from directional counts remains a separate task; the Weyl sector requires
additional nonlocal relational information.

\section{Model Status and Non-Claims}
\label{sec:status}

The following statements are established by the present implementation:
\begin{itemize}
  \item The conditional finite-window distribution on a fixed-count shell is
        exactly hypergeometric.
  \item The fall, structure, and promoted sectors form an exact partition.
  \item The recursive multi-clock ledger telescopes exactly.
  \item The matter and size curves are completely reproducible from the
        documented equations and default parameters.
  \item The classical and Dicke-state descriptions share the same
        fixed-count window marginal.
\end{itemize}

The present implementation does not establish:
\begin{itemize}
  \item an exact unconditioned solution of the one-flip discrete chain at
        every tick;
  \item a unique physical particle composition for any width;
  \item a derived value of the matter exponent $q$;
  \item a physical conversion from bit-flip coordinate to seconds or years;
  \item an FLRW metric, Hubble function, luminosity distance, or redshift law;
  \item the Einstein field equations or a quantized Ricci tensor;
  \item the observed dark-matter abundance or perturbation spectrum.
\end{itemize}

\section{Open Problems}
\label{sec:open}

\subsection*{Derive or eliminate the matter exponent}
The dominant immediate problem is equation~\eqref{eq:matter}.  A preferred
solution would derive matter from persistent worldline structure, so that
accidental equilibrium windows are not counted as particles and no external
power $q$ is required.

\subsection*{Derive widths and compositions}
The illustrative widths $(6,12,20)$ and the spectral default compositions
are placeholders.  They must be derived from the microstructure algebra and
its energy-information relation before particle names or abundance ratios
can be assigned.

\subsection*{Choose nested and parallel sectors physically}
A genuine formation chain should use the nested recursion.  Coexisting
species should use a joint or parallel allocation with an explicit rule that
prevents the same substrate bits from being assigned to incompatible
structures.  The current spectral code exposes both modes; the ontology must
select between them by sector.

\subsection*{Replace literal freeze by structural persistence}
Equation~\eqref{eq:survival} forbids every constituent bit from being touched.
A real stable structure may survive internal rearrangements.  The relevant
quantity is a first-passage or survival probability in the space of valid
composite states.

\subsection*{Implement the contrast and metric decoder}
The $n/2$ bound in equation~\eqref{eq:n2bound} should be connected to
operational rulers made from observer microstructures.  The result must
determine whether the global count represents a linear scale, an area, or a
volume and must define null propagation and redshift.

\subsection*{Derive the local Ricci sector}
The global subtraction in equation~\eqref{eq:size-code} must be promoted to a
local, directional conserved count and shown to have the continuum expansion
in equations~\eqref{eq:directional-count}--\eqref{eq:ricci-target}.

\subsection*{Complete the quantum measure}
The Dicke implementation supplies kinematics and exact marginal statistics.
A complete quantum cosmology also requires the amplitude or compression
measure over alternative static informational configurations and a derivation
of the relevant observation probabilities.

\section{Conclusion}
\label{sec:conclusion}

The current simulation implements a precise finite-budget model.  A
continuum Ehrenfest mean selects a rounded Hamming shell; the shell's
combinatorial entropy supplies the unfolded pool; exact conditional
hypergeometric statistics divide finite windows into falling,
structure-forming, and promoted sectors; and scale-dependent retarded clocks
organize those sectors into a nested hierarchy with an exact telescoping
ledger.  Matter is a weighted subset of the structure sector, and the
implemented size measure subtracts both matter bits and multi-clock backlog
from the unfolded pool.

The central physical suggestion is preserved without importing information
from outside the static universe: matter and spacetime are alternative
organizations of one finite informational budget.  Matter formation reduces
the free spatial-resolution proxy, while larger structures accumulate fewer
internal ticks per elementary transition.  The code therefore realizes the
proposed common origin of spatial contraction and slower clocks at the level
of a homogeneous statistical demonstrator.

The revision also makes the limitations explicit.  The late recovery of the
default curve depends on the non-derived $q=1$ matter weighting; the widths
and compositions are placeholders; the physical metric and proper-time map
have not been derived; and the model has not yet produced the Ricci tensor or
an observational cosmology.  These limitations identify concrete next
steps rather than undermining the exact combinatorial and accounting results
already present.

\appendix
\section{Code-to-Equation Map}
\label{app:code-map}

\begin{table}[ht]
\centering
\caption{Primary implementation functions and their mathematical roles.}
\small
\begin{tabular}{p{0.37\textwidth}p{0.55\textwidth}}
\toprule
Function or field & Mathematical content \\
\midrule
\texttt{p\_of\_tau} & Equation~\eqref{eq:p-code} with raw coordinate input. \\
\texttt{ones\_count\_from\_p} & Rounded mean shell, equation~\eqref{eq:k-code}. \\
\texttt{combinatorial\_entropy\_bits} & $\log_2\binom nk/n$. \\
\texttt{ground\_truth\_pool} & Returns $U(T)$ and $k(T)$. \\
\texttt{order\_parameter} & Equation~\eqref{eq:eta}. \\
\texttt{family\_fractions\_exact} & Equations~\eqref{eq:fall}--\eqref{eq:rise}. \\
\texttt{retard} & Equation~\eqref{eq:retard}. \\
\texttt{level\_state} & Equations~\eqref{eq:pending}--\eqref{eq:recursion}. \\
\texttt{matter\_content} & Equations~\eqref{eq:matter}--\eqref{eq:matter-total}. \\
\texttt{size\_measure} & Equation~\eqref{eq:size-code}. \\
\texttt{window\_marginal} & Equation~\eqref{eq:hypergeom}. \\
\texttt{class\_probability} & Equation~\eqref{eq:classprob}. \\
\texttt{pattern\_probability} & Equation~\eqref{eq:specificprob}. \\
\texttt{run\_cascade\_series} & Nested spectral composition hierarchy. \\
\texttt{run\_parallel\_series} & Independent coexisting spectral categories. \\
\bottomrule
\end{tabular}
\end{table}

\section{Reproduction Commands}
\label{app:commands}

The current classical figure is generated by
\begin{verbatim}
python simulations/cosmic_d.py \
  --n_bits 184 \
  --scales 6,12,20 \
  --matter_power 1 \
  --steps 3000 \
  --output emergent_structure_relativistic.png
\end{verbatim}
The spectral companion can be run as a nested cascade or with coexisting
parallel categories:
\begin{verbatim}
python simulations/cosmic_psi.py --scales 6,12,20 --output cascade.png
python simulations/cosmic_psi.py --parallel \
  --scales 6,12,20 --output parallel.png
\end{verbatim}
Specific spectral compositions may be supplied with, for example,
\begin{verbatim}
--compositions 1:5,2:10,3:17
\end{verbatim}
Such compositions remain modeling choices until independently derived.

\end{document}
